
\section{11. Conclusion}

\subsection{11.1 Summary}

We presented \textbf{NHA} (Native Hybrid Attention), a module that runs local 
sliding-window attention and global RFF attention in parallel, combined via 
a learned token-level gate.

\subsection{11.2 Results Recap}

\begin{table}[H]
\centering
\small
\begin{tabular}{llll}
\toprule
\textbf{Model} & \textbf{Acc $\Delta$} & \textbf{Complexity} & \textbf{Mean gate} \\
\midrule
DistilBERT & -0.46\% & $\mathcal{O}(L(w+m))$ & 0.75 \\
RoBERTa & -10.67\% & $\mathcal{O}(L(w+m))$ & 0.70 \\
GPT-2 & -0.58\% & $\mathcal{O}(L(w+m))$ & 0.72 \\
\bottomrule
\end{tabular}
\end{table}

\subsection{11.3 Key Takeaways}

1. \textbf{Linear complexity} in $L$ for both streams.
2. \textbf{Learned gates} outperform fixed fusion.
3. \textbf{Orthogonal RFF} stabilizes training.
4. \textbf{Chunked scan} preserves causality with bounded memory.
5. \textbf{No parameter reduction} -- compute-only optimization.

\subsection{11.4 Future Work}

\subsubsection{Short-Term}
\begin{itemize}
\item Learned window size
\item Per-head RFF count
\item Fused kernels for local + RFF
\item Long-context evaluation
\end{itemize}

\subsubsection{Long-Term}
\begin{itemize}
\item Theoretical error bounds
\item Cross-task gate transfer
\item Integration with quantization
\end{itemize}

\subsection{11.5 Closing Remarks}

NHA provides an interpretable hybrid attention mechanism for linear-time 
inference. We release code and checkpoints.

\subsection{11.6 Acknowledgments}

We thank {{ACKNOWLEDGMENTS}}.

\subsection{11.7 Author Contributions}

\begin{itemize}
\item \textbf{{{AUTHOR\_1}}} designed NHA, implemented the code, wrote the paper.
\item \textbf{{{AUTHOR\_2}}} provided feedback, edited.
\end{itemize}

\subsection{11.8 Conflict of Interest}

None.